\chapter*{Abstract}
\addcontentsline{toc}{chapter}{Abstract}

This study develops and evaluates OshiFit, a system that proposes ways of doing fan activities (\textit{oshi-katsu}) suited to each person, based on their personality traits and their current fan practices.

\paragraph{Background}
The number of people in Japan living with mental disorders has risen to about 4.19 million (2017), and hobbies are increasingly valued alongside sleep and exercise as a way of relieving stress. \textit{Oshi-katsu}, the sustained support of a favourite idol, character, or creator, is especially widespread among Generation Z. It gives fans a reason to live, a means of self-expression, and social connection, but it can also cause what fans call burnout: lost sleep, overspending, a sense of obligation to follow everything, and low mood from comparing themselves with other fans on social media. Screen-time limits and budgeting apps assume that less is better, yet in fan activities the practices that cause burden are often the same ones that make life meaningful, so blanket advice to cut down may remove the enjoyment as well. Moreover, fans rarely notice which of their own habits lead to burnout. Prior research has mostly described fans' motives and psychology; few studies have examined benefits and burdens together quantitatively, and none has supported individuals in adjusting how they engage. The author's previous study measured heart rate and stress during fan activities, but it involved only three participants over two weeks and covered mainly passive activities.

\paragraph{Aim and research questions}
This study analyses how personality traits, fan behaviours, well-being, and fan burnout relate across many fans, and builds a system that, based on these findings, offers proposals the user can choose and shows the predicted state after six weeks for keeping current habits (Plan A: current fan activities) and for adopting a proposal (Plan B: optimised fan activities), using a self-resembling avatar and scores with intervals. The aim is to help fans notice burdens they cannot easily see and to support self-chosen changes that raise well-being and reduce stress. Three research questions are addressed. RQ1: How are personality traits (TIPI-J) and fan behaviours associated with well-being and burnout, and how accurately can an individual's outlook be predicted? RQ2: Over six weeks of use, how do well-being and burnout change for fans who are already doing well (group A), who adopt Plan B (group B), and who keep Plan A (group C), within and between groups? RQ3: Do users recognise the sources of their burnout, and do those who changed their behaviour judge the system useful?

\paragraph{Survey (RQ1)}
An online survey of Japanese fans yielded 234 responses; after excluding 17 respondents with no fan activity and 2 who failed an attention check, 215 were analysed (113 men, 94 women, 8 other or undisclosed; median age 19). Personality (TIPI-J, 5-point), activity frequency, immersion, triggers for fan activity, well-being, and burnout were measured and analysed with k-means clustering, an observed-variable path model, moderation tests, and 10-fold cross-validated prediction. While 79\% said their \textit{oshi} gives them a reason to live, 40\% reported money worries and 33\% sleep disruption, and well-being and burnout were positively correlated ($r=.24$). Activity frequency was associated with both higher well-being ($\beta=.43$) and higher burnout ($\beta=.29$), as was immersion ($\beta=.20$, $.23$): benefits and burdens share the same sources. Personality was associated with how fans engage (extraversion with activity, $\beta=.19$; conscientiousness with lower immersion, $-.15$; neuroticism with mood-repair motives, $.21$; openness with weaker ones, $-.16$) but did not moderate the effect of activity ($p=.74$, $.79$). Clustering on personality alone separated poorly (silhouette 0.15--0.18). Cross-validated $R^2$ was .27 for well-being and .11 for burnout, showing that individual predictions are highly uncertain.

\paragraph{System}
Four design principles followed: (D1) never prescribe blanket reduction and always keep ``as is'' as a valid option; (D2) personalise by combining personality with the fan's own behaviours and burdens rather than assigning a fixed type; (D3) show benefits and burdens together; and (D4) be honest about uncertainty. From a 67-item questionnaire, OshiFit computes six axes on a 0--100 scale (well-being, enjoyment, sleep and daily life, money, obligation and relationships, social media and comparison) and shows each formula. Users with at least five axes $\geq 60$ and a WHO-5 raw score $\geq 13$ are told they are doing well (group A); others receive three of eleven proposals targeting their lowest axes, with ties broken by personality, and see Plan A and Plan B with 80\% intervals on a self-resembling avatar built on the device (the photo and landmarks are discarded). Users then choose Plan B (group B) or Plan A (group C) and complete weekly check-ins for six weeks.

\paragraph{Field study (RQ2, RQ3)}
After approval by the institutional ethics committee, a six-week field study is conducted. Groups are set by baseline state (A) and the user's own choice (B, C) rather than randomisation, because imposing changes on a hobby would contradict D1. The primary outcomes are WHO-5 and overall burnout, analysed within groups (baseline vs.\ week 6), across groups with A as reference (ANCOVA), between B and C with propensity-score weighting, and as a dose--response relationship within B. \textbf{[To be completed after the study: sample sizes, changes in WHO-5 and burnout, and ratings of awareness and usefulness.]}

\paragraph{Significance}
This study shows quantitatively that the benefits and burdens of fan activities arise from the same behaviours, so that ``do less'' support is unsuitable; provides evidence for using personality in combination with behaviour rather than as a fixed type; and builds a system that supports self-chosen behaviour change by comparing two futures on a self-resembling avatar while openly showing uncertainty. The approach may apply more broadly to hobbies and entertainment in which enjoyment and burden are intertwined, helping people sustain them rather than give them up.

\vspace{1em}
\noindent\textbf{Keywords:} oshi-katsu, fan activities, well-being, fan burnout, personality (TIPI-J), personal informatics, avatars, predictive uncertainty
